% ADASS 2026 proceedings, paper C301 (contributed talk, track 3).
% DRAFT of 2026-11-07. Based on the ADASS_template.tex LaTeX file, 24 Oct 2025,
% itself based on the ASP general author template (Copyright 2021, ASP
% Conference Series, revision 29 November 2021).
%
% Build with pdflatex, bibtex, pdflatex, pdflatex (or "make pdf", or Overleaf
% with the pdfLaTeX compiler and C301.tex as main document).
%
% Template rules kept in this draft:
%  - no package other than asp2023; no \renewcommand or \renewenvironment;
%  - references only as BibTeX entries in C301.bib, cited with \citet/\citep;
%  - figures are C301_f<n>.png in this folder; no non-ASCII characters;
%  - 4-page limit for contributed talks, references included.

\documentclass[11pt,twoside]{article}

% Do NOT use ANY packages other than asp2023.
\usepackage{asp2023}

\aspSuppressVolSlug
\resetcounters

\bibliographystyle{asp2021}

% Running head. More than 3 authors: "Surname1 et al."
\markboth{Lacopo et al.}{Multi-Agent Loop Engineering for Radio Imaging}

\begin{document}

\title{Multi-Agent Loop Engineering for Radio Astronomy Pipelines: Autonomous Refactoring and Scientific Validation}

% Five authors; Giovanni Lacopo is the first and presenting author.
% [PLACEHOLDER: affiliations of co-authors] Only the first author's affiliation
% is given in this draft. Before the final deadline add one \affil per
% institution (not repeated), the matching $^n$ superscripts after each name in
% \author, and complete the \paperauthor lines below. Do not guess institutions.
\author{Giovanni~Lacopo,$^1$ Emanuele~De~Rubeis, Claudio~Gheller, Giuliano~Taffoni, and Luca~Tornatore}
\affil{$^1$INAF -- Osservatorio Astronomico di Trieste, Via G.B. Tiepolo 11, 34143 Trieste, Italy; \email{giovanni.lacopo@inaf.it}}
% [PLACEHOLDER: affiliations of co-authors] \affil{$^2$...}

% One line per author; 9 arguments. [PLACEHOLDER: ORCID] third argument left
% blank for all authors. [PLACEHOLDER: co-author email, institution, department,
% city, state, postal code, country] left blank for the four co-authors.
\paperauthor{Giovanni~Lacopo}{giovanni.lacopo@inaf.it}{}{INAF -- Osservatorio Astronomico di Trieste}{}{Trieste}{}{34143}{Italy}
\paperauthor{Emanuele~De~Rubeis}{}{}{}{}{}{}{}{}
\paperauthor{Claudio~Gheller}{}{}{}{}{}{}{}{}
\paperauthor{Giuliano~Taffoni}{}{}{}{}{}{}{}{}
\paperauthor{Luca~Tornatore}{}{}{}{}{}{}{}{}

%\aindex{Lacopo,~G.}
%\aindex{De~Rubeis,~E.}
%\aindex{Gheller,~C.}
%\aindex{Taffoni,~G.}
%\aindex{Tornatore,~L.}


\begin{abstract}
We report the human-led refactoring of RICK, a distributed w-stacking dirty
imager built on MPI, OpenMP offload and HeFFTe, with large language model
(LLM) agents from two engines, Claude and OpenAI Codex. The researcher set
contracts and tolerances, ran every job on the Leonardo supercomputer and
judged the results; the agents implemented, reviewed and diagnosed. With one
writing agent, cross-engine review, no local execution and evidence-labelled
claims, the loop removed a
sign-inverting off-axis flux ramp and most of the radial droop with respect to
WSClean on a MeerKAT field, leaving a residual edge deficit, and produced a
streaming engine that splits work before it exceeds a known memory budget. On
VLA data of Abell~2255, CPU and GPU images differ in 34 of 81 million stored
pixels; agreement with WSClean is reported metric by metric.
\end{abstract}

%\ssindex{astronomy!radio!interferometry}
%\ssindex{algorithm!imaging}
%\ssindex{Artificial intelligence}
%\ssindex{software!development}
%\ssindex{software!testing}
%\ssindex{computing!high-performance computing (HPC)}
%\ssindex{computing!GPU}
%\ssindex{computing!OpenMP}
%\ssindex{libraries!MPI}
%\ssindex{observatories!ground-based!MeerKAT}
%\ssindex{observatories!ground-based!VLA}
%\ssindex{observatories!ground-based!SKA}

%\ooindex{WSClean, ascl:1408.023}
%\ooindex{casacore, ascl:1912.002}

\section{Introduction}

SKA-era imaging software must be correct, scalable and maintainable: a
refactor that is faster but two per cent wrong in flux is a scientific
failure. LLM agents can now modify real code bases \citep[e.g.][]{Jimenez2024},
but no unit test certifies a sky image. We asked whether agents can refactor
an HPC radio imager when every change is held to scientific gates set by a
human.

RICK \citep{DeRubeis2025,Lacopo2025,Lacopo2026} reads a MeasurementSet through
casacore \citep{2015A&C....12..174V}, grids weighted visibilities on w-planes
with a Kaiser--Bessel kernel \citep{Jackson1991}, transforms them with HeFFTe
\citep{Ayala2020}, corrects the w-term by w-stacking
\citep{2014MNRAS.444..606O}, and writes dirty images, PSFs and fitted beams to
FITS, combining MPI, OpenMP, an OpenMP-offload GPU build and optional NCCL
transport. The work, done in June--October 2026 on Leonardo \citep{Turisini2024},
covers dirty imaging only (no CLEAN or calibration); the reference is the
WSClean dirty image \citep{2014MNRAS.444..606O}.

\section{Method: Engineering the Loop}

Every change follows one loop. The researcher writes a contract: what must
stay invariant and which tolerance counts as a pass. One agent implements it,
the same engine and then the other review it statically, and only then does
the researcher run it on Leonardo. The agents never build or run RICK locally;
they diagnose from logs and FITS files copied back. Each iteration leaves a
design note and a commit stating its validation status, and every claim is
labelled \emph{code}, \emph{measurement} or \emph{unproven}.

The agents are versioned persona profiles: only an
\texttt{hpc-ska-specialist} (Claude and Codex profiles from one brief) writes
code; \texttt{cpp-hpc-reviewer}, \texttt{parallel-io-expert} and
\texttt{radioastronomy-expert} are read-only. One agent writes at a time:
parallel writers would produce conflicting edits in the gridding loops, where
a plausible but wrong change is hardest to spot. Three gates apply: agreement
with WSClean within 1\%, metric by metric; the same image for any MPI rank
count, so far enforced only on printed summaries (MeerKAT runs on 8, 12 and 32
ranks print the same peak and RMS); and execution choices such as chunking or
transport changing only floating-point evaluation order.

\section{Results}

\subsection{Reducing the Flux Gap on MeerKAT}

In earlier validation runs on the MeerKAT field PSZ2~G008.31 (4096$^2$ pixels
of 1$^{\prime\prime}$, uniform weighting), the peak ratio RICK/WSClean fell
off-axis, with a cosine ramp along $l-m$ that inverted the sign of 3 of 61
sources and a smooth radial droop. Five diagnoses, each re-checked against
WSClean, removed the ramp and most of the droop. The uniform-weighting density was binned per w-plane
instead of in 2D as in WSClean. The Kaiser--Bessel correction used the
analytic transform instead of inverting the sampled, truncated kernel. The
w-phase sign was wrong: code review found nothing, but replacing the w-phase
screen by the identity reproduced the single-plane image to floating-point
precision, clearing gridding, FFT and normalisation; the sign fix raised the
off-axis peak ratio from 0.35 to about 0.93. The ramp came from a half-cell
FFT offset shared by each visibility and its synthetic Hermitian conjugate,
giving a factor $\cos(\pi (l-m)/N_{\rm pad})$; a symmetric shift removed it
and all three sign inversions (Fig.~\ref{fig:meerkat}, left). Finally, 128
w-planes instead of 20, as in WSClean, reduced the droop coefficient 3.1-fold,
to a median ratio of 0.992 and a minimum of 0.967 (Fig.~\ref{fig:meerkat},
right). The minimum ratio of 0.967 is 3.3\% low, outside the 1\% gate; the
remaining deficit is under investigation. One fix often hid the next; we did
not record which engine proposed each.

\articlefiguretwo{C301_f1.png}{C301_f2.png}{fig:meerkat}{MeerKAT PSZ2~G008.31,
same-pixel peak ratio RICK/WSClean for sources above 50$\sigma$ (earlier
validation runs). \emph{Left:} against $l-m$; crosses and fitted cosine before
the Hermitian half-cell fix, filled circles after it. \emph{Right:} against
radius with 20 (light) and 128 (dark) w-planes, with quadratic fits.}

\subsection{Memory-Adaptive Streaming}

The most surprising development was not a speed-up: RICK now adapts to
the MeasurementSet instead of running out of memory. The contract is
\emph{count first, allocate second}: before allocating payload buffers, all
ranks exchange the numbers of samples they will read, send and receive and
compare them with every rank's known budget; if any rank would exceed it, all
ranks refuse the interval and split it. All stages are charged, including FFT
workspace and PSF, and an unknown budget stops the job at start-up. A static
Codex review challenged the first, Claude-written design before any run, and
its fixes shaped the final architecture. In the VLA GPU run below, one
interval of 16.6 million rows would have needed 80.1~GB on one GPU against a
57.5~GB ceiling; it was split, and the run completed with 16 intervals (4
planned). The 254~GB MeerKAT MeasurementSet also completed on 2 and 8 nodes.
This is a contract within our allocation model, not a proof against every
out-of-memory case: vendor FFT and NCCL allocations are only estimated.

\subsection{VLA Abell~2255: CPU, GPU and WSClean}

We imaged VLA data of Abell~2255 (circular feeds, 4500$^2$ pixels of
0.4$^{\prime\prime}$, 26 w-planes, Briggs robust $-0.5$ \citep{Briggs1995},
Stokes I, Q, U, V) on two nodes of each Leonardo partition
(Table~\ref{tab:vla}). With relative $L_2 = \|A-B\|_2/\|B\|_2$ over all
pixels, CPU and GPU images agree to $\leq1.54\times10^{-11}$ and the float64
PSFs to $\leq5.8\times10^{-15}$. The 2.22 time ratio comes from single runs on
different dates, with different toolchains, layouts, transports, interval
counts and timer scopes: it compares two configurations, not a general GPU
speed-up. Q, U and V agreement is a consistency check, not a polarimetric
validation.

\begin{table}[!ht]
\caption{VLA Abell~2255, single two-node runs. Times from RICK timers; pixel
metrics over all 20.25 million pixels of each float32 Stokes image.}
\label{tab:vla}
\smallskip
\begin{center}
{\small
\begin{tabular}{lcc}
\tableline
\noalign{\smallskip}
 & CPU: 32 ranks $\times$ 7 threads & GPU: 8 ranks, 8 A100 \\
\noalign{\smallskip}
\tableline
\noalign{\smallskip}
Total [s] & 1011.5 & 455.2 \\
Gridding, dirty + PSF [s] & 535.8 & 27.9 \\
Reading and preparation [s] & 261.8 & 319.5 \\
Unequal pixels, I/Q/U/V & \multicolumn{2}{c}{34 of $8.1\times10^{7}$ (30 by one float32 step)} \\
Max.\ $|$CPU$-$GPU$|$ [Jy\,beam$^{-1}$] & \multicolumn{2}{c}{$4.55\times10^{-13}$} \\
\noalign{\smallskip}
\tableline
\end{tabular}
}
\end{center}
\end{table}

Against the WSClean dirty image (Stokes I), both backends give the same
result, so the differences are not a CPU-versus-GPU effect. Within 1\%: the
peak, at the same pixel ($-0.10\%$); the robust noise $\sigma$ ($-0.17\%$); the
median ratio of the 4915 pixels above $10\sigma$ (1.0001); and 15 of the 24
local maxima above $50\sigma$. The other 9 deviate by 1.0--4.7\%, the largest
being the three sources beyond 2000 pixels from the centre
(Fig.~\ref{fig:vla}). The whole-image relative $L_2$ of 0.197 reflects coherent
rings and fringes in the residual map, not a flux error. The cause of the
edge droop remains open. Taken together, the 1\% gate is met by the global
metrics but not by all bright sources on either dataset (MeerKAT minimum ratio
0.967; 9 of 24 VLA maxima at 1.0--4.7\%), so the science gate is only
partially passed.

\articlefigure[width=0.5\textwidth]{C301_f3.png}{fig:vla}{VLA Abell~2255,
Stokes I: RICK/WSClean against radius. Line and band: median and interquartile
range per radial bin of the 4915 positive pixels above $10\sigma$, with
$\sigma = 1.4826\,{\rm MAD}$ of the WSClean image
(6.17~$\mu$Jy\,beam$^{-1}$). Dots: 24 local maxima above $50\sigma$ (same pixel). Shaded strip: $\pm1\%$.}

\section{Discussion and Conclusions}

Invariance tests (identity phase screen, rank counts, backends) found faults
that code review missed; a second-engine review of a design before any
GPU-hour was cheap insurance; and evidence labels told us how far to trust a
statement. Agents also produced plausible but wrong code, as the w-phase sign
shows, and the Leonardo round trip set the pace of the loop. Contracts,
tolerances, runs and the judgement of evidence stayed human. The evidence
comes from one project and one researcher, without a controlled comparison
against development without agents. Within these limits, LLM agents can
refactor a production HPC radio imager when held to scientific gates
(invariants, tolerances and a reference imager) rather than to code style, but
validation against WSClean on Leonardo remained essential for every result, and showed that the 1\% gate is not yet met everywhere. The researcher's role shifted from writing kernels to designing the
loop and judging its evidence.

% [PLACEHOLDER: ACKNOWLEDGEMENTS] e.g. Leonardo/CINECA computing allocation,
% funding, providers of the MeerKAT and VLA data, and any AI-use statement the
% editors may require. Co-authors are now listed in \author, so thank only
% people and institutions who are not authors. Keep the text on the same line
% as \acknowledgements, just before the bibliography. Do not invent grant
% numbers or project names.
%\acknowledgements [PLACEHOLDER]

\bibliography{C301}  % For BibTex

% if we have space left, we might add a conference photograph here. Leave commented for now.
% \bookpartphoto[width=1.0\textwidth]{foobar.eps}{FooBar Photo (Photo: Any Photographer)}

\end{document}
